\chapter{Ranges I --- Views, Range Adaptors, and Lazy Pipelines}
\label{ch:c20-ranges-i}

\begin{quote}
\emph{The ranges library is the second pillar and the one that most changes how you express data transformations. It replaces the iterator-pair calling convention of classic \texttt{<algorithm>} with composable, lazy pipelines built from the pipe operator. This chapter covers the range abstraction, what a view is and why it is cheap, the standard adaptors, and the lazy-evaluation model.}
\end{quote}

Ranges collapse ``copy\_if then a squaring loop'' into \texttt{data | views::filter(even) | views::transform(square)} --- a declarative pipeline that allocates nothing and computes each element on demand.

\section{The End of Iterator Pairs}
\label{sec:c20-end-iter-pairs}

\begin{lstlisting}[language=C++,caption={The old way vs. the ranges way}]
#include <vector>
#include <algorithm>
#include <ranges>
#include <iostream>

int main() {
    std::vector<int> nums{1, 2, 3, 4, 5, 6};

    // Old way: a temporary container and two passes
    std::vector<int> temp;
    std::copy_if(nums.begin(), nums.end(), std::back_inserter(temp),
                 [](int i){ return i % 2 == 0; });
    for (auto& x : temp) x = x * x;

    // C++20 ranges: one lazy pipeline, no temporary, single pass
    namespace views = std::views;
    auto result = nums
        | views::filter([](int i){ return i % 2 == 0; })
        | views::transform([](int i){ return i * i; });

    for (int i : result) std::cout << i << ' ';   // 4 16 36
}
\end{lstlisting}

\section{Ranges, Views, and the Core Concepts}
\label{sec:c20-ranges-views-concepts}

A \textbf{range} provides \texttt{begin()}/\texttt{end()}. A \textbf{view} is a range that is \textbf{cheap to copy and non-owning}.

\begin{center}
\begin{tabular}{llll}
\hline
\textbf{Term} & \textbf{Owns data?} & \textbf{Copy cost} & \textbf{Example} \\
\hline
Container (range) & yes & O(n) & \texttt{std::vector<int>} \\
View & no & O(1) & \texttt{views::filter(v, pred)} \\
\hline
\end{tabular}
\end{center}

Key concepts: \texttt{sized\_range}, \texttt{common\_range}, \texttt{viewable\_range}, and the iterator-strength ladder.

\section{The Pipe Operator and Range Adaptors}
\label{sec:c20-pipe-adaptors}

A \textbf{range adaptor} takes a range and produces a view. \texttt{range | adaptor} equals \texttt{adaptor(range)}.

\begin{lstlisting}[language=C++,caption={Function-call form and pipe form are equivalent}]
#include <ranges>
#include <vector>

std::vector<int> v{1,2,3,4,5};
namespace vws = std::views;

auto a = vws::transform(vws::filter(v, [](int x){ return x>2; }),
                        [](int x){ return x*10; });

auto b = v | vws::filter([](int x){ return x>2; })
           | vws::transform([](int x){ return x*10; });
\end{lstlisting}

\section{The Standard Views Catalogue}
\label{sec:c20-views-catalogue}

Several popular views (\texttt{zip}, \texttt{enumerate}, \texttt{chunk}, \texttt{slide}, \texttt{join\_with}) are \textbf{C++23, not available here}.

\begin{center}
\begin{tabular}{ll}
\hline
\textbf{View} & \textbf{Effect} \\
\hline
\texttt{all} & view over an entire range \\
\texttt{filter(pred)} / \texttt{transform(f)} & select / map \\
\texttt{take(n)} / \texttt{take\_while(pred)} & leading elements \\
\texttt{drop(n)} / \texttt{drop\_while(pred)} & skip leading elements \\
\texttt{reverse} & reversed traversal \\
\texttt{join} & flatten range-of-ranges \\
\texttt{split(delim)} & tokenize on delimiter \\
\texttt{elements<N>} / \texttt{keys} / \texttt{values} & project tuple fields \\
\texttt{iota(a[,b])} & arithmetic sequence \\
\texttt{single(x)} / \texttt{empty<T>} & degenerate generators \\
\texttt{counted(it,n)} / \texttt{common} & misc adaptors \\
\hline
\end{tabular}
\end{center}

\begin{lstlisting}[language=C++,caption={Composing several adaptors}]
#include <ranges>
#include <vector>
#include <iostream>

int main() {
    std::vector<int> v{5, 1, 8, 2, 9, 3, 7};
    namespace vws = std::views;
    auto pipeline = v
        | vws::filter([](int x){ return x > 2; })
        | vws::transform([](int x){ return x - 1; })
        | vws::take(3)
        | vws::reverse;
    for (int x : pipeline) std::cout << x << ' ';     // 8 7 4
}
\end{lstlisting}

\section{Lazy Evaluation: What Actually Happens During Iteration}
\label{sec:c20-laziness}

Building a pipeline does no work; computation happens element-by-element \emph{during iteration}.

\begin{lstlisting}[language=C++,caption={Laziness: each element flows through the whole chain on demand}]
#include <ranges>
#include <vector>
#include <iostream>

int main() {
    std::vector<int> v{1, 2, 3, 4};
    namespace vws = std::views;
    auto pipe = v
        | vws::filter([](int x){ std::cout << "filter " << x << '\n'; return x % 2 == 0; })
        | vws::transform([](int x){ std::cout << "transform " << x << '\n'; return x * x; });
    for (int x : pipe) std::cout << "got " << x << "\n";
}
\end{lstlisting}

Infinite views are usable when bounded downstream; the optimizer fuses the chain into one loop. C++20 ships views (lazy), not eager actions.

\section{Generating Views: iota and Friends}
\label{sec:c20-iota}

\begin{lstlisting}[language=C++,caption={iota generates sequences; bounded and infinite}]
#include <ranges>
#include <iostream>

int main() {
    namespace vws = std::views;
    for (int i : vws::iota(0, 5)) std::cout << i << ' ';   // 0 1 2 3 4
    std::cout << '\n';
    auto first_squares = vws::iota(1)
                       | vws::transform([](int n){ return n*n; })
                       | vws::take(5);                       // 1 4 9 16 25
    for (int s : first_squares) std::cout << s << ' ';
}
\end{lstlisting}

\section{Decomposing and Splitting: join, split, elements}
\label{sec:c20-join-split}

\begin{lstlisting}[language=C++,caption={join flattens, split tokenizes, keys/values project}]
#include <ranges>
#include <vector>
#include <string>
#include <map>
#include <iostream>

int main() {
    namespace vws = std::views;
    std::vector<std::vector<int>> nested{{1,2},{3,4},{5}};
    for (int x : nested | vws::join) std::cout << x << ' ';   // 1 2 3 4 5

    std::string text = "a,bb,ccc";
    for (auto tok : text | vws::split(','))
        std::cout << std::string_view(tok.begin(), tok.end()) << '|';

    std::map<std::string,int> m{{"x",1},{"y",2}};
    for (auto& k : m | vws::keys)   std::cout << k << ' ';
    for (auto& val : m | vws::values) std::cout << val << ' ';
}
\end{lstlisting}

\texttt{views::split} tokens are not always contiguous; build \texttt{string\_view} from \texttt{tok.begin()/end()}.

\section{Performance and the Materialization Question}
\label{sec:c20-materialization}

\begin{itemize}
\item \textbf{No intermediate allocation} is the headline win.
\item \textbf{Re-traversal cost}: \texttt{filter} recomputes its predicate each pass; \texttt{begin()} is O(n). Materialize if iterated repeatedly.
\item \textbf{No \texttt{std::ranges::to} in C++20} (that is C++23).
\end{itemize}

\begin{lstlisting}[language=C++,caption={Materializing a view into a container in C++20}]
#include <ranges>
#include <vector>

auto view = std::views::iota(1, 6) | std::views::transform([](int n){ return n*n; });

std::vector<int> out1(view.begin(), view.end());    // requires common_range
std::vector<int> out2;
for (int x : view) out2.push_back(x);
// C++23 would allow: view | std::ranges::to<std::vector>(); -- NOT C++20
\end{lstlisting}

\section{Professional Insights}
\label{sec:c20-ranges-i-insights}

\textbf{Default to pipelines for clarity, verify codegen for hot paths.} Laziness plus inlining usually matches a hand-written loop; confirm once with the disassembler in a latency-critical loop.

\textbf{Materialize a filtered view you iterate more than once.} \texttt{filter} re-evaluates its predicate and pays O(n) for \texttt{begin()} each traversal.

\textbf{Mind the C++20 materialization gap.} \texttt{std::ranges::to} is C++23; materialize via iterator-pair constructor or a loop.

\textbf{Know which views are not in C++20.} \texttt{zip}, \texttt{enumerate}, \texttt{chunk}, \texttt{slide}, \texttt{join\_with} are C++23.

\textbf{Exploit laziness deliberately with infinite generators.} \texttt{views::iota(0)} plus \texttt{take} is the idiomatic allocation-free replacement for index loops.
